\section{Conclusion and Future Work}
\balance

The inability to access comprehensive medical records significantly hinders the development of Deep Learning in healthcare. In this paper, we successfully bridged this gap by developing a robust, rule-based generation engine that produced 500 000 highly realistic, multimodal synthetic patient records simulating an IoMT triage environment. By benchmarking an MLP and a Tabular Transformer on this dataset, we demonstrated that our synthetic data is complex, sparse, and exhibits realistic overlapping clinical phenotypes. Furthermore, the application of SHAP visually validated that the generated data strictly adheres to physiological coherence.

However, a notable limitation of our current framework is the absence of unstructured medical imaging data and longitudinal behavioral metrics. In modern clinical settings, radiological imaging, such as Chest X-rays, CT scans, or MRIs, is often the definitive modality for diagnosing conditions like Pulmonary Embolism or Hemorrhagic Stroke. Furthermore, patient lifestyle factors, including sleep patterns, daily physical activity, and consumption habits, play a crucial role in predicting disease onset and assessing baseline health.

Therefore, future work will focus on expanding our synthetic generation pipeline across two distinct axes. First, we plan to include paired medical images, potentially leveraging advanced Generative AI (e.g., Diffusion Models), and integrating Computer Vision models (such as CNNs or Vision Transformers) to extract features from these scans. Second, we aim to incorporate continuous lifestyle and behavioral data, mimicking the longitudinal tracking capabilities of consumer-grade IoMT wearables. Concatenating these high-dimensional spatial and behavioral modalities with our existing acute tabular framework will pave the way for a comprehensive multimodal diagnostic assistant.